\bookinterlude{IV}{V}{O bilhete no meio da mesa}
Seis meses depois, chega um pedido de reposição. A cliente escreve “Faça outra igual à que mandei de volta para polir”. Pego o cartão que deixei ao lado do telefone. Tem o torno que está livre, a última medida que entregamos e o custo de um dia de atraso. Até ontem eu usava esse cartão para decidir se valia a pena pagar hora extra. Hoje, em que medida você vai regular o torno?

Abro a pasta da cliente. A primeira peça media 5,2 mm, a segunda, 5,8. O carimbo do polimento está na ficha da segunda. Faço uma cópia da pasta e mudo o carimbo de ficha. Na cópia, o torno que está livre é o mesmo, a última entrega continua sendo a de 5,8 mm e agora a cliente tem razão de esperar uma peça de 5,2.

Dá vontade de acrescentar “última peça polida” ao cartão do telefone. Faço isso. Na manhã seguinte chega outro pedido da mesma cliente. Desta vez ela quer uma peça igual à primeira que aprovou depois de mudar de endereço. Volto a abrir a pasta.

No primeiro dia havia espaço de sobra para duas moedas nessa mesa. Hoje eu preciso do espaço para espalhar as fichas. Ponho o bilhete novo no meio delas. Você lê duas linhas, procura um pouco e põe o dedo sobre uma ficha. Que palavras no bilhete levaram você até ela?
